% =====================================================================
%  Harald Høffding, H. Martensen: Jacob Böhme (review)
%  In: Ude og Hjemme.  Nordisk illustreret Ugeblad, udgivet af Otto Borchsenius og F. Hendriksen.  Femte Aargang 1881-1882 (Kjøbenhavn: F. Hendriksen; J. Jørgensen & Co.), Nr. 238 [23 April 1882], pp. 365-366.
%  Antiqua, two columns.
%
%  TRANSCRIPTION -- GENERATED; DO NOT EDIT BY HAND.  Høffding's three pieces of 1879-82 (Filosofien i Sverig; Om
%  Pessimismen; the review of Martensen's Jacob Böhme) are built by ONE pipeline (ebook-method/hoeffding-1879-82/) into a
%  combined working file, which split.py cuts into three transcriptions.
%  Rebuild from the repo root:   sh ebook-method/hoeffding-1879-82/rebuild.sh
%
%  STATUS: COMPLETE (2026-10-15): the volume's title page and the whole text, pp. 365-366.
%
%  EDITION.  The only printing.  Text of record: Internet Archive ude-og-hjemme-aarg-5 (a private scan, 300 ppi JP2, uploaded 2024);
%  extract ~/bibliotek/Høffding, Harald/1881-82-ude-og-hjemme-5-nr238-pp365-366.pdf (SCANS.tsv, sha256), made by extract.py (title, Indholdsfortegnelse,
%  pp. 365-366).  The volume's Indholdsfortegnelse: „Høffding, Harald: ... H. Martensens »Jacob Bøhme« .... 365“.  The
%  date of Nr. 238 is inferred (Nr. 237 is dated Søndagen den 16. April 1882; the paper appeared weekly).  The review's
%  head runs across both columns of p. 365 below the end of a story; the review ends on p. 366 above „Det kongelige Theater“.
%
%  ACCURACY (2026-10-15; one collation for all three texts, 68 printed pages, ~28,100 words).
%  (a) No second copy of any of the three was found (Internet Archive, Nasjonalbiblioteket, Project Runeberg, Google Books;
%      HathiTrust not reachable).  The DRAFT is our own tesseract reading (tessdata_best, 300 dpi; dan+swe for Nordisk
%      tidskrift, dan for the others), chosen by measurement (words unattested in the corpus, and for Nordisk tidskrift
%      a majority of three readings: measure.py), laid over the working PDF and turned into LaTeX by draft.py.  The
%      WITNESS, collated word by word (collate/check.py): NB's ALTO (ABBYY) for Nordisk tidskrift and Vor Ungdom, the
%      Archive's tesseract for Ude og Hjemme.  1,213 disagreements in all.
%  (b) Two tesseract readings that share faults: the text was therefore PROOFREAD WHOLE at the image, column by column,
%      against the draft (reader P1: 26 corrections, both planted errors found; proofapply.py).  Quotation marks »...« as
%      printed.  Structure (head, ending, signature) read by reader K; letterspacing by P1 and K.
%  PAGE CHECK.  SAMPLED PAGE CHECK (2026-10-15): 4 pages drawn at random from four strata (Filosofien i Sverig pp. 35, 60; Om Pessimismen
%      p. 83; Ude og Hjemme p. 366, left column), 1,595 words transcribed blind at the image by a reader who saw no other
%      text (pagecheck.py): 1 word error (ɔ: read as o:, then mended in all five places of the class), 1 punctuation error
%      („Følelse“. read Følelse,), 2 italic runs (a posteriori, a priori) -- all mended; the reader's 2 other differences were his
%      own misreadings at a zoom.  About 2 errors per 1,000 words before mending.
%  PRINTED AS IS (misprints kept, marked at the line):
%    p. 365: „levendé“ for levende, an accent printed; zoom z-marks2.png
%  CONVENTIONS.  Page turns: % --- p. N --- and \opage{N} (the printed page); the columns are marked % (p. N, column c); a word divided over a page turn is
%  split with %.  Footnotes as \footnote at their marks, numbered per page as printed (no notes).  \emph{} =
%  letterspacing, \textit{} = italics.  The running heads and the sheet signatures are not transcribed.
% =====================================================================
\documentclass[12pt,a4paper,openany]{book}
\usepackage[utf8]{inputenc}
\usepackage[T1]{fontenc}
\usepackage{libertinus}
\usepackage{libertinust1math}
\usepackage[danish]{babel}
\usepackage{textalpha}   % Greek typed directly
\usepackage[protrusion=true,expansion=true]{microtype}
\usepackage{setspace}
\usepackage[margin=1.3in]{geometry}
\usepackage{fancyhdr}
\usepackage{hyperref}
\hypersetup{pdftitle={Anmeldelse af H. Martensen, Jacob Böhme}, pdfauthor={Harald Høffding}, hidelinks}
\pagestyle{fancy}
\fancyhf{}
\fancyhead[LE,RO]{\thepage}
\fancyhead[C]{\textit{H. Høffding: Anmeldelse af H. Martensen, Jacob Böhme (1882)}}
\fancypagestyle{plain}{\fancyhf{}\fancyhead[LE,RO]{\thepage}}
\setstretch{1.3}
\usepackage{marginnote}
\usepackage{graphicx}
\DeclareUnicodeCharacter{0254}{\reflectbox{c}}   % ɔ (id est)
\makeatletter

\makeatother
\newcommand{\opage}[1]{\marginnote{\footnotesize\textit{[#1]}}}
\newcommand{\signature}[1]{\par\smallskip\hfill #1\hspace*{2em}\par}
\pdfstringdefDisableCommands{\renewcommand{\footnote}[2][]{}\renewcommand{\opage}[1]{}}

\begin{document}

\begin{titlepage}
\begin{center}
% [a ruled frame round the page]
{\Huge UDE OG HJEMME.}\\[0.5em]
{\large NORDISK ILLUSTRERET UGEBLAD.}\\[1em]
% [an engraved cherub with a garland]
{\scriptsize UDGIVET AF}\\[0.3em]
{\sffamily OTTO BORCHSENIUS {\small OG} F. HENDRIKSEN.}\\[0.6em]
% [wavy rule]
{\small\bfseries FEMTE AARGANG 1881---1882.}\\
% [wavy rule]
\vfill
{\small KJØBENHAVN.}\\[0.3em]
{\scriptsize TRYKT PAA F. HENDRIKSENS FORLAG HOS J. JØRGENSEN \& CO.}
\end{center}
\end{titlepage}
\clearpage

% --- p. 365 ---
\setcounter{footnote}{0}%
\opage{365}\begin{center}{\large H. \emph{Martensen: Jacob Böhme.}}\\[0.3em]{\small Theosofiske Studier.  Kjøbenhavn.  Gyldendalske Boghandels Forlag.  1881.}\end{center}\medskip

% (p. 365, column 1)
\setcounter{footnote}{0}%
Den protestantiske Kirkes Grundlæggelse var et af de
betydningsfuldeste Resultater af den almindelige Higen efter
Frigjørelse, Inderlighed og Naturlighed, som udmærkede
den Tid, man har kaldet Gjenfødelsens (Renaissancens) Tid.
Men næppe havde den unge Kirke sikret sin Existens, før
den mente at besidde Frihed og Inderlighed nok. Den
formulerede sine Dogmer og tilbagetrængte enhver fri
Tanke- og Følelsesretning. Skjønt Luther selv var dybt
paavirket af den mystisk-spekulative Retning, som udfoldede
sig i Tyskland fra det 14de Aarhundrede af, og hvis mest
fremragende Tænker var Mester Eckhart, saa lede de
Mystikere, der opstod indenfor den nye Kirke, dog en
haard Medfart. Schwenckfeld og Sebastian Franck førte et
omflakkende og ustadigt Liv for deres Anskuelsers Skyld,
og Valentin Weigel levede kun uantastet, fordi han ikke
lod sine Skrifter trykke. Paa Universiteterne herskede
Orthodoxien og den af Melanchton grundede udvortes
Kombination af Theologi og aristotelisk Filosofi. Men de
mystiske Bøger og Manuskripter gik fra Haand til Haand
mellem Lægmænd. Sandsynligvis er den mærkelige Skomager
Jacob Böhme fra Görlitz (f. 1575, d. 1624) bleven
indviet i disse Skrifter, medens han rejste paa Haandværket.
Vandrende Haandværkssvende hørte til de vigtigste
Bærere og Udbredere af de mystiske Ideer og Skrifter.
Med den religiøse Mystik forbandt sig allerede hos Weigel
Indflydelsen af Paracelsus' Alkymi og Naturfilosofi, som et
realistisk Gjæringsstof, der blev tilsat den mystiske Idealisme.
Denne dobbelte Bevægelse førtes til sin højeste og betydningsfuldeste
Form, da den slog ned i Böhmes dybe, barnlige
Sind. Religiøs Grublen, storartet Fantasi, levendé Natursans, % PRINTED AS IS (levendé: for levende, an accent printed; zoom z-marks2.png)
psykologisk Blik og skarp Tænkning i Forening med
en vidunderlig Originalitet og Kraft i Sproget gjorde ham
til en enestaaende Skikkelse i Tænkningens Historie. Han
er den mythologiske Forgænger for den tyske Filosofi, hvis
Grundtanker han udtaler i mystisk Form.

% (p. 365, column 2)
\setcounter{footnote}{0}%
Det er denne Mand, hvis Liv og Lære danner Hovedindholdet
af det foreliggende Skrift. Der er et aandeligt
Slægtskab mellem Forfatteren og hans Helt. Martensens
Tænkning har stedse havt et afgjort mystisk eller theosofisk
Præg. Naar i Theologien Ordet saa ofte maa gjøre Fyldest
for det manglende Begreb, saa er det hos Martensen vistnok
oftest det prægnante Billede, den anelsesfulde Fantasi, der
danner Tankens Afslutning og tilslører, hvad den ikke
formaar at udfylde. Hvad der stiller ham saa højt over
hele den yngre Slægt af Theologer, er, at han dog endnu
har Mod til at tænke. Han bebrejder de yngre Theologer,
at de blot holde sig til den overleverede Kirkelære og ikke
indlade sig paa at vinde en sammenhængende Erkjendelse
og Forstaaelse. Det er smukt at se den gamle Forkæmper
for theologisk Spekulation blive sin Ungdoms Kjærlighed
tro til det Sidste. Ligesom et af de Arbejder, med hvilke
han begyndte sin glimrende Forfattervirksomhed, var et
Skrift om Mester Eckhart, saaledes har han nu i sin høje
Alderdom sat sig ved Fødderne af den dybsindige Skomager
fra Görlitz for at skildre den afsluttende Repræsentant
for hin mystisk-theosofiske Bevægelse og for selv at
spinde Traaden videre. Thi Skriftet er ikke en rent
historisk Monografi; Forf. kriticerer ikke blot sin Helt,
men søger ogsaa at gaa videre i de af ham anviste Spor.

Hvor Biskoppen og Skomageren ere uenige, er jeg for
mit Vedkommende mest tilbøjelig til at tage Skomagerens
Parti; og hvor de ere enige, bliver jeg ofte uenig med
dem begge. Men det forekommer mig dog, at man læser
det foreliggende Skrift med ulige større Interesse og Udbytte
end Forf.s »Kristelige Ethik«. I »Ethiken« færdedes
man i den virkelige Verden og saa' den tilspidsede højkirkelige
Tankegang tørne mod alle Nutidens Frihedsbestræbelser.
Her bevæger man sig derimod saa højt over
Virkeligheden, at man kun en sjælden Gang øjner Jorden
som et forsvindende Atom. Vi føres ind i Udviklings%
% --- p. 366 ---
\setcounter{footnote}{0}%
\opage{366}processer, Kriser og Katastrofer i den oversanselige Verden,
skildrede med saadant Liv og Fynd, at naar vi komme til
det rigtige, menneskelige Syndefald, forekommer det os at
være en meget ubetydelig Historie, en Storm i et Glas
Vand. Men i hin Verden er Forf. som i sit Element.
Hans Talent er udpræget dogmatisk-theosofisk, er et Talent
for billedskuende Tænken, en Kombination af Følelse,
Fantasi og Forestilling, der egner sig bedre til en Kamp
med Sproget for at udtrykke det Utænkelige og Usigelige
end til indtrængende Sympathi for og Forstaaelse af reale,
skarpt bestemte Virkelighedsforhold.

Det er det store Problem om det Ondes Oprindelse,
der fører Böhme og med ham Martensen op i hin oversanselige
Verden. Ved at se Kampen mellem det Gode og
det Onde i Naturen blev Böhme greben af dybt Tungsind
og henfaldt til Tvivl og Grublen. Hvorledes kan denne
Kamp forenes med Guds Godhed og Enhed? Böhme fandt
Løsningen i den allerede af Mester Eckhart udtalte Tanke,
at uden Modsætning er ingen Selvbevidsthed mulig, eller
som Böhme paa sit kjærnefulde Sprog udtrykker sig: »Kein
Ding ohne Widerwertigkeit mag ihm selber offenbar werden.«
Skal Gud altsaa være den klart bevidste og villende Gud,
maa der være en Modsætning at overvinde, et Mørke, der
kan danne Lysets Baggrund, en Kval, der kan lade Glæden
aabenbare sig, et Nej, der kan fremhæve Ja'et. I denne
Naturgrund indenfor Gud finder han »Helvedes Fundament«,
Spiren til det Onde. Ved Lucifers Fald blev det frigjort
og optændt til en Ild, der brænder i Alt, hvad der er
skadeligt og ondt i Naturen og i Mennesket. Hele denne
Udviklingsproces skildres af Forf. efter Böhme paa en
storartet og betagende Maade. Man læser Fremstillingen
som et Digterværk, et protestantisk Sidestykke til Dantes
guddommelige Komedie (kun at det her er en Tragedie).

Böhme er en dristig Tænker, trods al sin Barnefromhed
og sin lutherske Orthodoxi. Han gyser ikke tilbage
for nogen Tanke, der ligger paa hans Vej. Selv den
bibelske Fremstilling forkaster han, naar den »wider die
Philosophia und Vernunft laufet«. Den mosaiske Skabelseshistorie
læmpes efter den Paracelsiske Naturfilosofi, og det
Ondes Spire forfølges tilbage til Guddommens inderste
Natur. Her nedlægger Martensen sin Protest. Böhme er
ham for naturalistisk og panteistisk; Böhme tænker »efter
en hedensk, panteistisk, kosmogonisk Type«. Martensen
fordrer, at Gud skal tænkes som \emph{virkelig fra først af},
Virkelighed forud for al Mulighed. Jeg sér ikke Andet,
end at denne Fordring gjør al Theosofi umulig. Theosofien
spørger dér, hvor de Fleste forlængst have hørt op
at spørge; Guddommens Virkelighed forudsætter for Theosofien
lige saa vel en (i Principet, om ikke i Tiden) forud
gaaende Mulighed, som al anden Virkelighed gjør. Vi
% (p. 366, column 2)
\setcounter{footnote}{0}%
forstaa kun Virkeligheden gjennem de Muligheder, af hvilke
den udfolder sig. Skal altsaa Gudsbegrebet tænkes, kan
Muligheden ikke holdes ude. Mulighed og Virkelighed ere
af den Slags Begreber, der ikke kunne skilles ad og anvendes
hver for sig. Det er lige saa vilkaarligt at begynde
med en absolut Virkelighed som med en absolut Mulighed.
Vi kjende kun Virkelighed i Forhold til Mulighed, og
omvendt. Böhme har Ret i, at \emph{ville} vi tænke, maa vi
anvende Mulighedsbegrebet, statuere en Evolution fra
Mulighed til Virkelighed. Men han har Uret i, at vi kunne
naa til nogen Afslutning her: thi ligesom enhver Virkelighed
forudsætter en Mulighed, forudsætter enhver Mulighed en
given Virkelighed. --- Martensen kalder Böhmes Lære for
hedensk. Böhme har foregrebet et saadant Angreb i
følgende Udtalelse (som jeg ikke har fundet i den foreliggende
Bog, skjønt den er karakteristisk): »Nun wirst du
sagen, ich schreibe heidnisch. Höre und siehe und merke
den Unterschied wie dies alles sei; denn \emph{ich schreibe nicht heidnisch sondern philosophisch}.« --- Historisk
sét er Beskyldningen for Hedenskab vistnok urigtig.
Thi det er netop hedenske Tænkere (Platon og Aristoteles),
der mest energisk have forsøgt at tænke Guddommen som
evig Virkelighed; og i Kristendommens Lære om en lidende
Gud kommer den modsatte Tanke frem.

For at undgaa al Panteisme hævder Martensen, at
Gud ikke skabte Verden, fordi det var en Nødvendighed
for ham. Netop derfor lægger han saa stor Vægt paa
Spekulationerne om en oversanselig Natur, for at Skabelsen
af den sanselige Verden kan staa som noget saa Vilkaarligt
som muligt. Det synes, som om han her bestandig
arbejder sig selv imod; thi er Kjærlighed virkelig Motivet
til Skabelsen, saa maa denne ogsaa være fremgaaet af en
Trang. Den Kjærlighed, der kan undvære sit Værk og er
lige glad, om det bliver til Noget eller ikke, give vi her i
denne »sanselige« Verden ikke Meget for. Maaske den
højere Logik, som Theosofien raader over, korrigerer vore
Forestillinger om Kjærlighed (ligesom om Virkelighed).
De, der ere glade ved nogenlunde at kunne tænke efter
den »endelige«, lavere Logik, maa da give Afkald paa at
forstaa hine høje Ting.

Böhmes blivende Betydning beror ikke mindst paa den
psykologiske Skarpsindighed, hvormed han har opdaget og
fastholdt Modsætningens, Kontrastens, Relativitetens Betydning
i al Erfaring, i alt aandeligt Liv og Bevidsthed. Hans
System er en genial, som Theologi forklædt Psykologi. Det
afgiver et indirekte Bevis for, at vore Begreber ikke egne
sig til Løsning af overnaturlige Gaader. Theosofien maa
først grundlægge en ny Psykologi og en ny Logik, før
den kan naa sit Maal.
\signature{\textbf{\sffamily Harald Høffding.}}

\end{document}
